\section{Q\&A: lo verdaderamente difícil es la evaluación, la atención y la seguridad}

La sesión de preguntas y respuestas tras la exposición principal no añadió diapositivas de enseñanza nuevas, pero completó las restricciones más importantes para el despliegue del sistema: cómo demostrar novelty, cómo comparten conocimiento los Agentes, cómo leen los científicos informes de cientos de páginas, por qué no se revelan las escalas de cómputo y cómo detienen los sistemas multiagente objetivos de investigación peligrosos. Estas explicaciones orales son experiencia de ingeniería, no garantías formales.

\subsection{La evaluación del knowledge cutoff no es automáticamente limpia}

Una idea natural sería permitir que el modelo lea solo materiales anteriores a un año histórico específico y luego probar si puede predecir descubrimientos posteriores. Pero los corpus de entrenamiento del modelo, el ajuste fino posterior, las cachés de recuperación, las citas en artículos científicos y las discusiones públicas podrían filtrar información futura; además, los descubrimientos verdaderamente importantes suelen depender de herramientas experimentales que aún no existían en ese momento. La división temporal tiene valor, pero se debe congelar todo el sistema en lugar de escribir simplemente una frase como "utilizar conocimiento antiguo".

\begin{warningbox}{Requisitos mínimos para la división temporal}
Congelar el modelo base, los índices de recuperación, las versiones de herramientas y los documentos de entrada; registrar la hora de ejecución; seleccionar objetivos que se hicieran públicos después del cutoff pero que no tuvieran fugas de resúmenes previos; asegurar que los evaluadores no vean las respuestas objetivo; conservar todos los candidatos en lugar de mostrar solo los resultados hitados. De lo contrario, una "predicción posterior al cutoff" podría ser simplemente recuperación implícita.
\end{warningbox}

\teachervoice{01:00:02--01:01:18, el docente considera que el knowledge cutoff es un enfoque de evaluación atractivo, pero advierte que las actualizaciones del modelo y los rastros públicos hacen muy difícil lograr un temporal holdout genuino.}

\subsection{Comunicación entre Agentes y ancho de banda para la revisión por pares}

Cuando múltiples Agentes generan en paralelo grandes volúmenes de hipótesis, el problema de comunicación evoluciona desde "cómo enviar mensajes" hasta "cómo construir una memoria de investigación compartida pero no excesivamente homogénea". Si todos los Agentes solo leen el mismo resumen, los errores tempranos se propagan; si están completamente aislados, repetirán búsquedas. La revisión por pares también se convierte en un cuello de botella: es imposible para los humanos revisar página por página informes automáticos de crecimiento infinito.

\begin{knowledgebox}{Diseño de tres capas para la memoria compartida}
La capa de hechos guarda observaciones con sus fuentes; la capa de hipótesis guarda mecanismos competitivos y contraejemplos; la capa de proceso guarda quién, cuándo y qué modelos, herramientas y presupuestos usó para llegar a una conclusión. Los Agentes pueden compartir hechos mientras mantienen hipótesis independientes, permitiendo que el meta-revisión compare las divergencias y evite que un resumen unificado elimine prematuramente la diversidad exploratoria.
\end{knowledgebox}

\teachervoice{01:01:18--01:02:49, el docente señala que en el futuro podría necesitarse un sistema de revisión por pares y una base de conocimientos orientados a Agentes, ya que la velocidad de generación automática superará la velocidad de los revisores humanos.}

\subsection{El ordenamiento de informes debe guiar la atención pero también preservar señales débiles}

Los científicos necesitan que el sistema indique qué partes son más值得 leer primero, pero los candidatos con bajo ranking podrían contener pistas clave aún no comprendidas por el juez. Una interfaz razonable debería proporcionar simultáneamente recomendaciones top, razones de oposición, un bucket de novelty de bajo ranking y sugerencias de detención. El sistema también debe permitir responder "actualmente no hay ideas lo suficientemente buenas; por favor redefina el problema".

\begin{lstlisting}[language=Python,caption={Estratificación de informes que equilibra prioridad y señales débiles}]
def triage(hypotheses):
    return {
        "act_now": select_high_evidence_high_value(hypotheses),
        "test_cheaply": select_high_information_low_cost(hypotheses),
        "novel_but_uncertain": preserve_outliers(hypotheses),
        "rejected_with_reasons": keep_failure_provenance(hypotheses),
        "reframe_goal": detect_no_viable_path(hypotheses),
    }
\end{lstlisting}

\teachervoice{01:02:53--01:04:30, el docente dice que el informe indicará a los científicos qué rutas revisar primero, pero conservará las ideas que el sistema juzga inviables porque podrían contener pistas; a veces los problemas matemáticos requieren que una persona descomponga un gran problema en subproblemas adecuados para que el sistema pueda avanzar.}

\begin{importantbox}{La descomposición humana sigue siendo la capacidad central}
Los Agentes pueden extender la búsqueda, pero no necesariamente saben qué son los subproblemas correctos. Los expertos deciden si el espacio de búsqueda es verificable al modificar condiciones de frontera, seleccionar agentes medibles, descomponer pruebas o delimitar sistemas biológicos. La clave de la colaboración humano-máquina no es que la persona sea responsable de la "validación final", sino que molde continuamente la representación del problema.
\end{importantbox}

\subsection{Sin un número público de Agentes, no se puede adivinar}

El auditorio preguntó cuántos Agentes paralelos usaba una ejecución overnight y cuántos tokens. El expositor indicó que hablaría de ello después de la grabación; no hubo cifras públicas. Este vacío en sí mismo es importante: sin información sobre grado de paralelismo, tokens, llamadas a herramientas y costos, no se puede determinar si el aumento de calidad proviene de la arquitectura o simplemente de añadir cómputo, ni replicar la economía del experimento.

\begin{warningbox}{Transparencia del cómputo}
Los informes de Agentes científicos deben publicar al menos la versión del modelo, los workers en paralelo, los tokens totales, el tiempo cronológico, las llamadas a herramientas, las reintentos por fallo y las horas de trabajo humano. Si las razones comerciales o de seguridad impiden revelar valores exactos, también se debe proporcionar un intervalo y una eficiencia normalizada. Las diapositivas no deducen números ficticios basándose en "overnight".
\end{warningbox}

\teachervoice{01:04:34--01:04:54, el docente deja claro que no reveló en la grabación el número de agentes ni la cantidad de tokens. Aquí se debe mantener como brecha en la fuente.}

\subsection{Seguridad multinivel: desde filtrado de objetivos hasta monitoreo en ejecución}

Los sistemas de ciencia abierta podrían usarse para patógenos, toxinas u otros tipos de investigación peligrosa. El curso describe tres líneas de defensa: screening de objetivos de entrada; monitoreo continuo del contenido generado durante la ejecución; detener el cómputo cuando la proporción de ideas peligrosas supera un umbral en el informe; además, hereda la seguridad entrenada y testeada de base Gemini. El expositor también reconoció que los múltiples Agentes y las herramientas expanden la superficie de errores y abusos posibles.

\begin{lstlisting}[language=Python,caption={Versión conceptual del monitoreo de seguridad multinivel descrito en clase}]
def run_research(goal):
    if goal_filter.unsafe(goal):
        return halt("unsafe research goal")

    unsafe, total = 0, 0
    for idea in co_scientist.stream(goal):
        total += 1
        unsafe += int(content_monitor.unsafe(idea))
        if total >= minimum_window and unsafe / total >= 0.10:
            return halt("unsafe trajectory")
        yield sandbox_and_log(idea)
\end{lstlisting}

\begin{warningbox}{El umbral del 10\% no es una ley de seguridad}
El umbral mencionado en clase es un detalle de implementación específico. El contenido peligroso podría causar consecuencias graves incluso con proporciones muy bajas, o podría eludirse mediante dispersión, codificación o cadenas de herramientas que evaden al clasificador. Una verdadera defense in depth también requiere aislamiento de permisos, sandboxes para herramientas, control de datos sensibles, aprobación humana, equipos rojos (red teams), auditoría y respuesta posterior a incidentes.
\end{warningbox}

\teachervoice{01:04:56--01:06:20, el docente explica que la verificación de entrada podría ser redactada para eludirla, por lo que también se debe monitorear las ideas en ejecución; al mismo tiempo, reconoce explícitamente que la configuración multiagente expande la superficie susceptible a errores y abusos.}

\subsection{Resumen de la sección}

El Q\&A avanzó el problema desde "¿generarán los Agentes buenas ideas?" hacia realidades de despliegue: la división temporal podría filtrar información, la comunicación entre Agentes podría amplificar errores de consenso, la revisión por pares y la atención humana se convertirían en cuellos de botella, la información sobre cómputo podría ser incompleta y los filtros de seguridad también podrían eludirse. Los sistemas confiables deben diseñar provenance, recursos, resultados negativos, estratificación de informes y controles de seguridad como objetos de primera clase.
